%Chargement des paquets
\usepackage{amsmath}
\usepackage{amsfonts}
\usepackage{amssymb}
\usepackage{enumerate}
\usepackage{mathtools}
\usepackage{bbm}
\usepackage{xparse, etoolbox}
\usepackage{enumerate}
\usepackage{enumitem}
\usepackage{mathabx}
\usepackage{babel}[french]

%environment
\usepackage{amsthm}
\usepackage{keytheorems}
\usepackage{hyperref} 

%Interface théorème
\newkeytheorem{lem}[
	name = Lemme
]

\newkeytheorem{prop}[
	name = Proposition
]

\newkeytheorem{thm}[
	name = Théorème
]

\newkeytheorem{coro}[
	name = Corollaire
]

\renewcommand*{\proofname}{Démonstration}

\theoremstyle{definition}
\newtheorem{defn}{Définition}

\theoremstyle{definition}
\newtheorem*{nota}{Notation}

\theoremstyle{definition}
\newtheorem*{limi}{Liminaire}

\theoremstyle{remark}
\newtheorem{rmq}{Remarque}

\theoremstyle{plain}
\newtheorem*{fait}{Fait}


%Ensembles de nombres
\newcommand{\N} {\mathbb{N}}
\newcommand{\Ne}{\N^\ast}
\newcommand{\Z} {\mathbb{Z}}
\newcommand{\Q} {\mathbb{Q}}
\newcommand{\R} {\mathbb{R}}
\newcommand{\Rb}{\overline{\mathbb{R}}}
\newcommand{\Rp}{\R_+}
\newcommand{\Rm}{\R_-}
\newcommand{\K} {\mathbb{K}}
\newcommand{\Ket} {\mathbb{K^{\ast}}}
\newcommand{\Cx}{\mathbb{C}}

%Ensembles, tribus
\newcommand{\A}{\mathbb{A}}
\renewcommand{\P}{\mathcal{P}}
\newcommand{\T}{\mathcal{T}}
\newcommand{\F}{\mathcal{F}}
\newcommand{\C} {\mathcal{C}}
\newcommand{\ouv}{\mathcal{O}}
\newcommand{\B}{\mathcal{B}}
\newcommand{\M}{\mathcal{M}}
\newcommand{\etag}{\mathcal{E}}
\newcommand{\etagOp}{\etag^{+}(\Omega)}
\newcommand{\etagO}{\etag(\Omega)}
%%Lp avant quotientage
\newcommand{\Lic}{\mathcal{L}^1}
\newcommand{\Lpc}{\mathcal{L}^p}
\newcommand{\Linfc}{\mathcal{L}^{\infty}}
\newcommand{\Lifull}{\Lic(\Omega, \B, m)}
\newcommand{\Lpfull}{\Lpc(\Omega, \B, m)}
\newcommand{\Linffull}{\Linfc(\Omega, \B, m)}
%%Lp après quotientage
\newcommand{\Li}{L^1}
\newcommand{\Ld}{L^2}
\newcommand{\Lp}{L^p}
\newcommand{\Linf}{L^{\infty}}
\newcommand{\Listd}{\Li(\Omega, m)} 
\newcommand{\Ldstd}{\Ld(\Omega, m)} 
\newcommand{\Lpstd}{\Lp(\Omega, m)} 

%Quelques opérateurs
\DeclareMathOperator{\codim}{codim}
\DeclareMathOperator{\rg}{rg}
\DeclareMathOperator{\tr}{tr}
\DeclareMathOperator{\vect}{Vect}
\DeclareMathOperator{\ima}{Im}
\DeclareMathOperator{\id}{Id}
\DeclareMathOperator{\sign}{sign}
\DeclareMathOperator{\ext}{ext}
\newcommand{\ind}{\mathbbm{1}}
\newcommand*{\spr}[2]{\left(\,#1\,\middle|\,#2\,\right)}
\newcommand{\interior}[1]{%
  {\kern0pt#1}^{\mathrm{o}}%
}
\DeclareMathOperator{\card}{card}
\DeclareMathOperator{\ppcm}{ppcm}

%Commandes de pure flemme
\newcommand{\veps}{\varepsilon}
\newcommand{\intab}{\int_{a}^{b}}
\newcommand{\intR}{\int_{-\infty}^{\infty}}
\newcommand{\intinf}{\int_{- \infty}^{+ \infty}}
\newcommand{\equi}{\Leftrightarrow}
\newcommand{\Kev}{$\K$-espace vectoriel }
\newcommand{\Kevs}{$\K$-espaces vectoriels }
\newcommand{\Rev}{$\R$-espace vectoriel }
\newcommand{\Revs}{$\R$-espaces vectoriels }
\newcommand{\Cev}{$\Cx$-espace vectoriel }
\newcommand{\Cevs}{$\Cx$-espaces vectoriels }
\newcommand{\evn}{(E, \norm{.})}
\newcommand{\interf}[2]{[#1,#2]}
\newcommand{\limb}{\lim\limits_}
\newcommand{\lebn}{\lambda_n}
\newcommand{\pp}{\text{~presque partout}}
\newcommand{\derivp}[2]{\frac{\partial #1}{\partial #2}}
\newcommand{\famiu}{(u_i)_{i \in I}}
\newcommand{\sev}{\text{~s.e.v.}}
\newcommand{\Mnp}{M_{n,p}(\K)}
\newcommand{\Mn}{M_{n}(\K)}
\newcommand{\tq}{\text{~t.q.~}}
\newcommand{\mq}{~montrer~que~}
\newcommand{\et}{\quad \text{et} \quad}
\newcommand{\ou}{\text{~ou~}}
\newcommand{\Kx}{\K[X]}


%Commandes de gars chelou
\renewcommand{\subset}{\subseteq}

%Commande restriction, ty duckduckgo
\def\restr#1#2{\mathchoice
              {\setbox1\hbox{${\displaystyle #1}_{\scriptstyle #2}$}
              \restraux{#1}{#2}}
              {\setbox1\hbox{${\textstyle #1}_{\scriptstyle #2}$}
              \restraux{#1}{#2}}
              {\setbox1\hbox{${\scriptstyle #1}_{\scriptscriptstyle #2}$}
              \restraux{#1}{#2}}
              {\setbox1\hbox{${\scriptscriptstyle #1}_{\scriptscriptstyle #2}$}
              \restraux{#1}{#2}}}
\def\restraux#1#2{{#1\,\smash{\vrule height .8\ht1 depth .85\dp1}}_{\,#2}} 

%Quelques normes
\DeclarePairedDelimiter\abs{\lvert}{\rvert}
\DeclarePairedDelimiter\labs{\bigl\lvert}{\bigr\rvert}
\DeclarePairedDelimiter\norm{\lVert}{\rVert}
\DeclarePairedDelimiterXPP\normi[1]{}\lVert\rVert{_1}{\ifblank{#1}{\:\cdot\:}{#1}}
\DeclarePairedDelimiterXPP\normd[1]{}\lVert\rVert{_2}{\ifblank{#1}{\:\cdot\:}{#1}}
\DeclarePairedDelimiterXPP\normp[1]{}\lVert\rVert{_p}{\ifblank{#1}{\:\cdot\:}{#1}}
\DeclarePairedDelimiterXPP\normq[1]{}\lVert\rVert{_q}{\ifblank{#1}{\:\cdot\:}{#1}}
\DeclarePairedDelimiterXPP\norminf[1]{}\lVert\rVert{_\infty}{\ifblank{#1}{\:\cdot\:}{#1}}

%Bracket en angle
\DeclarePairedDelimiter\angl{\langle}{\rangle}

%Set, parenthèses
\newcommand{\set}[1]{\{ #1 \}}
\newcommand{\bigset}[1]{\bigl\{ #1 \bigr\}}
\newcommand{\Bigset}[1]{\Bigl\{ #1 \Bigr\}}
\newcommand{\bigpar}[1]{\bigl( #1 \bigr)}
\newcommand{\Bigpar}[1]{\Bigl( #1 \Bigr)}

%Opitmisation des opérateurs de comparaison
\DeclarePairedDelimiter\tio{o (}{)}
\DeclarePairedDelimiter\gro{O (}{)}


\pagestyle{fancy}

\fancyhead[C]{Dénombrement des polynômes unitaires irréductibles sur un corps fini}
\fancyhead[R]{}

\newcommand{\D}{\mathcal{D}}
\newcommand{\U}{\mathcal{U}}
\newcommand{\Fp}{\mathbb{F}_p}
\newcommand{\Fpd}{\mathbb{F}_{p^d}}

\begin{document}

\underline{Source :} Rombaldi - Algèbre - pages 422 à 424 (331 pour Möbius et 375 pour la structure d'algèbre). \newline

\underline{Leçons :}
\begin{itemize}
\item    123 : Corps finis. Applications.
\item    125 : Extensions de corps. Exemples et applications
\item    141 : Polynômes irréductibles à une indéterminée. Corps de rupture. Exemples et applications.
\item    190 : Méthodes combinatoires, problèmes de dénombrement.
\end{itemize}

\begin{nota}
Pour tout nombre premier $p \geq 2$, on note $\Fp$ le corps $\Z/p\Z$. 
Pour tout entier $n \geq 1$, on note $\U_n(p)$ l'ensemble de tous les polynômes unitaires irréductibles de degré $n$ dans $\Fp[X]$,
et $I_n(p)$ le cardinal de cet ensemble. On notera aussi $\D_n$ l'ensemble des diviseurs de $n$ dans $\Ne$.
Enfin, on note $P_n(X) = X^{p^n} - x \in \Fp[X]$.
\end{nota}

\section{Résultats préliminaires}

La démonstration de ces résultats peut éventuellement être passée à l'oral.

\subsection{Fonction de Möbius}

\begin{defn}[fonction de Möbius]
On définit la fonction de Möbius $\mu$ sur $\Ne$ par :
\[ \left \{ \begin{array}{c c} 
1 & \text{ si } n = 1 \\
(-1)^r & \text{ si } n = \prod_{i=1}^r p_i \\
0 & \text{ sinon}.
\end{array} \right. \]
Cette fonction est donc une suite réelle définie sur $\Ne$.
\end{defn}

\begin{defn}[produit de convolution]
Le produit de convolution de deux suites réelles $u, v$ définies sur $\Ne$ est la suite $u \ast v$ définie par :
\[ \forall n \in \Ne, \quad (u \ast v)(n) = \sum_{d \in \D_n} u(d) v \left ( \dfrac{n}{d} \right) . \]
On note $e$ l'élément neutre pour le produit de convolution, on a $e=(1, 0, \dots )$.
\end{defn}

\begin{lem}
Le produit de convolution est commutatif.
\end{lem}

\begin{proof}
Soient $u, v$ deux suites réelles définies sur $\Ne$, et $n \in \Ne$, on a :
\[ (u \ast v)(n) = \sum_{d \in \D_n} u(d) \left ( \dfrac{n}{d} \right) = \sum_{d' \in \D_n} u \left ( \dfrac{n}{d'} \right)  = (v \ast u)(n) . \]
\end{proof}

\begin{lem}
En notant $\omega$ la suite constante égale à 1, on a :
$\mu \ast \omega = e$.
\end{lem}

\begin{proof}
Pour $n=1$, on a $(\mu \ast \omega)(1) = \mu(1) \omega(1) = 1$ 
et pour $n \geq 2, (\mu \ast \omega)(n) = \sum_{d \in \D_n} \mu(d)$. \\ 
Notons $n = \prod_{i=1}^r p_i^{\alpha_i}$ la décomposition en facteurs premiers de $n$, 
tous les diviseurs $d$ de $n$ sont de la forme :
\[ d = \prod_{i=1}^r p_i^{\beta_i} , \]
avec chaque $\beta_i$ compris entre 0 et $\alpha_i$. 
Or $\mu(d) = 0$ si un $\beta_i$ est strictement plus grand que $1$, on déduit que :
\[ (\mu \ast \omega)(n) = \sum_{(\beta_1, \dots, \beta_r) \in \set{0,1}^r} \mu \left ( \prod_{i=1}^r p_i^{\beta_i} \right ) . \]
On s'intéresse donc aux $r$-uplets $(\beta_1, \dots, \beta_r)$ formés de 0 et de 1. 
On sait qu'il y a exactement $\binom{r}{k}$ $r$-uplet comportant exactement $k$ 1,
et pour un tel $r$-uplet la fonction de Möbius prend la valeur $(-1)^k$. 
On en déduit que :
\[ \forall n \geq 2, \quad (\mu \ast \omega)(n) = \sum_{k=0}^r \binom{r}{k} (-1)^k = (1-1)^r = 0 . \]
soit $\mu \ast \omega = e$.
\end{proof}

\begin{thm}[formule d'inversion de Möbius]
\label{thm:invMob}
Soient $u, v$ deux suites réelles, les assertions ci-dessous sont équivalentes :
\begin{align}
    \forall n \in \Ne, \quad u(n) = \sum_{d \in \D_n} v(d) \label{eq:m1} \\
    \forall n \in \Ne, \quad v(n) = \sum_{d \in \D_n} \mu(d) u \left (\dfrac{n}{d} \right ) \label{eq:m2}
\end{align}
\end{thm}

\begin{proof}
Supposons que l'assertion (\ref{eq:m1}) soit vraie. On a alors $u = v \ast \omega$, ce qui implique que 
$u \ast \mu = v \ast \omega \ast \mu = v$ soit $v = \mu \ast u$, par commutativité, soit l'assertion (\ref{eq:m2}).
On démontre la réciproque de façon identique.
\end{proof}

\section{Dénombrement des polynômes}

\begin{lem}
\label{lem:divis}
Soit $p \geq 2$ un nombre premier et $n \geq 1$ un entier. 
Alors $p$ divise tous les coefficients binomiaux $\binom{p^n}{k}$, avec $k \in \set{0, \dots, p^n-1}$.
\end{lem}

\begin{proof}
Pour $k \in \set{0, \dots, p-1}$, on a :
\[ \binom{p}{k} = \dfrac{p \cdot (p-1) \cdots (p-k+1)}{k \cdot (k-1) \cdots 1} = \dfrac{a}{k!} \in \N , \]
On peut écrire $a = \binom{p^n}{k} k!$.
On remarque que $p$ et $k!$ sont premiers entre eux, or $p$ divise $a$, donc (théorème de Gauss), $p$ divise $\binom{p^n}{k}$. \\
Considérons maintenant le polynôme $(1+X)^{p^n}$. 
Pour $n=1$, on déduit de ce que l'on vient de démontrer que 
\begin{equation} (1+X)^p \equiv 1 + X^p \pmod p , \label{eq:3} \end{equation}
Supposons qu'on ait, pour $n \geq 1$ quelconque :
\begin{equation} (1+X)^{p^n} \equiv 1 + X^{p^n} \pmod p , \label{eq:4} \end{equation}
alors, on écrit :
\begin{align*}
(1+X)^{p^{n+1}} & = \left ( (1+X)^{p^n} \right )^p \equiv (1 + X^{p^n})^p \pmod p \qquad \text{(d'après (\ref{eq:4}))} \\
 & \equiv 1 + X^{p^{n+1}} \pmod p \qquad \text{(d'après (\ref{eq:3}))} 
\end{align*}
On a donc démontré par récurrence que $(1+X)^{p^n} \equiv 1 + X^{p^n} \pmod p$ pour tout $n \geq 1$.
En particulier, cela signifie que tous les coefficients binomiaux $\binom{p^n}{k}$ sont divisibles par $p$.
\end{proof}

\begin{thm}
\label{thm:structpoly}
Soit $\K$ un corps et $P \in \K[X]$ un polynôme unitaire de degré $n \geq 1$. 
Alors $\K[X]/(P)$ est une algèbre de dimension $n$ et $(\overline{X^k})_{0 \leq k \leq n-1}$ en est une base. 
De plus, si le polynôme $P$ est irréductible, alors $\K[X]/(P)$ est un corps.
\end{thm}

\begin{proof}
Le fait que $\K[X]/(P)$ soit un anneau résulte de la compatibilité des lois de l'anneau $\K[X]$ avec les classes d'équivalence modulo $P$.
On vérifie facilement que cet anneau est commutatif et unitaire. 
De plus, on peut le munir d'une opération externe en définissant :
\[ \forall (\lambda, \overline{A}) \in \K \times \K[X]/(P), \quad \lambda \overline{A} = \overline{\lambda A} . \]
Ensuite, tout polynôme de $A \in \K[X]$ s'écrit $A = PQ + R$ avec $R$ de degré au plus $n-1$. 
On déduit donc que $\overline{A} = \overline{R} \in \vect(\overline{X^k})_{0 \leq k \leq n-1}$, la famille exhibée est génératrice.
Supposons qu'il existe une combinaison linéaire telle que :
\[ \sum_{k=0}^{n-1} a_k \overline{X^k} = \overline{0} , \]
alors le polynôme $a_{n1} X^{n-1} + \dots + a_0$ est multiple de $P$, donc, par contrainte de degré, est nul, 
donc tous les $a_k$ sont nuls et on déduit que la famille est libre donc une base. \newline

Suppons que $P$ soit irréductible et prenons $A \in \K[X]$ non nul. On note $R$ le reste de la division euclidienne de $A$ par $P$.
Comme $P$ est irréductible et de degré strictement supérieur à $R$, ces deux polynômes sont premiers entre eux et, d'après le théorème 
de Bézout, il existe $U, V \in \K[X]$ tels que $RU + PV = 1$. 
Comme $\overline{A} = \overline{R}$, on a $\overline{A}\overline{U} = \overline{1}$. 
En considérant que $\K[X]$ est appliqué surjectivement dans $\K[X]/(P)$, 
on a bien démontré que tout élément de ce dernier admet un inverse, c'est donc un corps.
\end{proof}


\begin{lem}
\label{lem:divPn}
Tout diviseur irréductible de $P_n$ dans $\Fp[X]$ est de degré divisant $n$. Réciproquement, pour tout diviseur $d$ de $n$, 
tout polynôme $P \in \U_d(p)$ divise $P_n$.
\end{lem}

\begin{proof}
Pour tout $P \in \U_d(p)$, selon le théorème $\ref{thm:structpoly}$, l'ensemble $\Fpd = \Fp[X]/(P)$ est un corps 
ainsi qu'un espace vectoriel de dimension $d$. C'est donc un corps fini de cardinal $p^d$. \newline

Soit $P$ un diviseur irréductible de $P_n$, on note $d$ le degré de $P$. A une constante multiplicative près, $P \in \U_d(p)$,
on peut donc ajouter cette hypothèse sans perte de généralité. 
Selon la remarque ci-dessous, $\Fpd$ est un corps et $\overline{P_n} = \overline{0}$ dans ce corps, autrement dit, 
$\overline{X}^{p^n} = \overline{X}$. 
Considérons $\overline{Q} \in \Fpd$, on peut l'écrire :
\[ \overline{Q} = \sum_{k=0}^{d-1} a_k \overline{X}^k , \]
avec les $a_k$ des éléments de $F_p$. D'après le lemme $\ref{lem:divis}$, on a :
\[ \overline{Q}^{p^n} = \left ( \sum_{k=0}^{d-1} a_k \overline{X}^k \right )^{p^n} 
= \sum_{k=0}^{d-1} a_k^{p^n} \left ( \overline{X}^k \right )^{p^n} , \]
avec de plus $a_k^{p^n} = a_k$ dans $\Fp$ d'après le théorème de Fermat, et :
\[ \left ( \overline{X}^{p^n} \right )^k = \overline{X}^k , \]
ce qui donne :
\[ \overline{Q}^{p^n} = \sum_{k=0}^{d-1} a_k \overline{X}^k = \overline{Q} . \]
Ainsi, pour tout polynôme $Q$ de $\Fp[X]$ non multiple de $P$ :
\[ \overline{Q}^{p^n-1} = \overline{1} . \]
$\Fpd^{\ast}$ est un groupe multiplicatif d'ordre $p^d-1$, donc pour tout élément $\overline{Q} \in \Fpd^{\ast}$ on a 
$\overline{Q}^{p^d-1} = \overline{1}$. Autrement dit, $p^n-1$ est un multiple de $p^d-1$.
Supposons que $d$ ne divise pas $n$, on écrit $n=qd + r$ avec $0<r<d$ et :
\[ p^n-1 = p^qd p^r -1 = p^r ( p^{qd} - 1) + p^r - 1 , \]
or $p^d - 1$ divise $p^n -1$ et $p^{qd}-1$, donc va aussi diviser $p^r-1$ ce qui est impossible ($p^r-1 < p^d-1$).
Finalement, $d$ divise $n$ ce qui achève la démonstration de la condition nécessaire.

Réciproquement, soit $d$ un diviseur de $n$ et soit $P \in \U_d(p)$. 
$\Fpd^{\ast}$ est un groupe multiplicatif d'ordre $p^d-1$,  donc, d'après le théorème de Lagrange :
\[ \forall X \in \Fp[X], \quad \overline{X}^{p^d-1} = \overline{1} . \]
On en déduit que $\overline{X}^{p^d} = \overline{X}$, ce qui permet de démontrer (par récurrence) que :
\[ \forall k \in \Ne, \quad \overline{X}^{p^kd} = \overline{X} . \]
Comme $d$ divise $n$, on a, en particulier $\overline{X}^{p^n} = \overline{X}$. 
Autrement dit, le polynôme $P$ divise $X^{p^n} - x = P_n(X)$.
\end{proof}

\begin{thm}
Le polynôme $P_n(X) = X^{p^n} - X$ est sans facteur carré dans $\Fp[X]$ et on a la décomposition en facteurs irréductibles :
\[ X^{p^n} - X = \prod_{d \in \D_n} \prod_{P \in \U_d(p)} P . \]
\end{thm}

\begin{proof}
Comme $P'_n(X) = p^n X^{p^n-1} - 1 = - 1$ dans $\Fp[X]$, les polynômes $P_n$ et $P'_n$ sont premiers entre eux et $P_n$ est sans 
facteurs carrés dans $\Fp[X]$. 
La formule exhibée dans le théorème est donc une conséquence directe du lemme $\ref{lem:divPn}$.
\end{proof}

\begin{thm}
Pour tout $n \in \Ne$, on a : 
\[ n I_n(p) = \sum_{d \in \D_n} \mu \left ( \dfrac{n}{d} \right ) p^d . \]
\end{thm}

\begin{proof}
En prenant les degrés dans la formule du théorème précédent on trouve :
\[ \forall n \Ne, \quad p^n = \sum_{d \in \D_n} \sum_{0 \in \U_d(p)} \deg(P) = \sum_{d \in \D_n} d I_d(p) , \]
En utilisant la formule d'inversion de Möbius (cf. théorème $\ref{thm:invMob}$), on peut écrire :
\[ \forall n \Ne, \quad n I_n(p) = \sum_{d \in \D_n}  \mu(d) p^{n/d} = \sum_{d \in \D_n} \mu \left ( \dfrac{n}{d} \right ) p^d . \]
\end{proof}

\begin{coro}
Pour tout entier naturel $n$ non nul, il existe dans $\Fp[X]$ des polynômes irréductibles de degré $n$ et on a :
\[ I_n(p) \sim_{n \to \infty} \dfrac{p^n}{n} . \]
\end{coro}

\begin{proof}
On veut tout d'abord montrer que $I_n(p) > 0$ pour tout entier $n$ non nul. \\
Pour $n=1$, on a $I_1(p) = p \geq 2$. Pour $n=2$, on a $I_2(p) = \dfrac{p(p-1)}2 \geq 1$. \\
Pour $n \geq 3$, on a 
\[ n I_n(p) = p^n + \sum_{d \in \D_n \setminus \set{n}} \mu \left ( \dfrac{n}{d} \right ) p^d . \]
Or :
\begin{align*}
\abs*{  \sum_{d \in \D_n \setminus \set{n}} \mu \left ( \dfrac{n}{d} \right ) p^d } 
& \leq \sum_{d \in \D_n \setminus \set{n}} \abs{*\mu \left ( \dfrac{n}{d} \right )} p^d \\
& \leq \sum_{d \in \D_n \setminus \set{n}} p^d \qquad (\mu \text{ à valeurs dans } \set{-1,0,1}) \\
& \leq \sum_{k=1} ^{E[n/2]} p^k \qquad (*) \\
& \leq p \dfrac{p^{E[n/2]} - 1}{p-1} < p^{E[n/2]} - 1 \leq p^{n/2+1} .
\end{align*}
L'inégalité (*) étant obtenue en considérant qu'un diviseur de $n$ strictement inférieur à $n$ est, au plus égal à $n/2$. 
Finalement, on peut minorer $n I_n(p)$ par :
\[ n I_n(p) > p^n - p^{n/2+1} > 0 . \]
Enfin, en considérant que
\[ \abs*{ \dfrac{n I_n(p)}{p^n} -1 } < \dfrac{p^{n/2+1}}{p^n} \to_{n \to \infty} 0 , \]
on a bien $I_n(p) \sim_{n \to \infty} \dfrac{p^n}{n}$.
\end{proof}

\end{document}